\subsection{与正则序列相伴的扩张类}

设 A 是一个环，M 是一个 A-模， \(x = (x_{1}, \ldots, x_{n})\) 是 A 中元素的序列。记 \(M_{i}\) 为 A-模 \(\mathbf{M}/(x_{1}\mathbf{M} + \cdots + x_{i-1}\mathbf{M})\) ，其中 \(i = 0, \ldots, n + 1\) ，使得 \(M_{0}\) 与 \(M_{1}\) 都与 M 等同，且 \(\mathbf{M}_{n+1} = \mathbf{M}/(\mathbf{x})\mathbf{M}\) 。设

\[
\overline {{x}} _ {i}: \mathbf {M} _ {i - 1} \to \mathbf {M} _ {i}, \qquad i = 1, \dots , n,
\]

表示由 \(M_{i-1}\) 上比率为 \(x_{i}\) 的位似与 \(M_{i-1}\) 到 \(M_{i}\) 上的典范投影复合而成的 A-同态。最后，我们记 \(p: M_{n} \to M/(x)M\) 为典范投影。图表

\[
0 \longrightarrow \mathrm{M} \xrightarrow {\overline {{{x}}} _ {1}} \mathrm{M} _ {1} \xrightarrow {\overline {{{x}}} _ {2}} \mathrm{M} _ {2} \longrightarrow \dots \xrightarrow {\overline {{{x}}} _ {n}} \mathrm{M} _ {n} \xrightarrow {p} \mathrm{M} / (\boldsymbol {x}) \mathrm{M} \longrightarrow 0\tag{19}
\]

是一个正合列当且仅当序列 x 是 M-正则的。此后假设序列 x 对 M 是正则的。与正合列 (19) 相伴的元素 \(\theta_{x} \in \operatorname{Ext}_{A}^{n}(\mathbf{M}/(\mathbf{x})\mathbf{M}, \mathbf{M})\) 也称为与 M-正则序列 x 相伴。

设 i 为整数， \(1 \leqslant i \leqslant n\) 。注意，序列 (19) 可以分解为两个正合列

\begin{enumerate}
\def\labelenumi{(\arabic{enumi})}
\setcounter{enumi}{19}
\tightlist
\item
\end{enumerate}

\[
0 \longrightarrow \mathrm{M} \xrightarrow {\overline {{x}} _ {1}} \mathrm{M} _ {1} \xrightarrow {\overline {{x}} _ {2}} \mathrm{M} _ {2} \longrightarrow \dots \xrightarrow {\overline {{x}} _ {i}} \mathrm{M} _ {i} \longrightarrow \mathrm{M} / (x _ {1} \mathrm{M} + \dots + x _ {i} \mathrm{M}) \longrightarrow 0\tag{21}
\]

\[
0 \longrightarrow \mathrm{M} / (x _ {1} \mathrm{M} + \dots + x _ {i} \mathrm{M}) \xrightarrow {\overline {{x}} _ {i + 1}} \mathrm{M} _ {i + 1} \longrightarrow \dots
\]

\[
\longrightarrow \mathrm{M} _ {n} \stackrel {p} {\longrightarrow} \mathrm{M} / (x) \mathrm{M} \longrightarrow 0
\]

它们正是与序列 \((x_{1}, \ldots, x_{i})\) （它对 M 是正则的）以及与序列 \((x_{i+1}, \ldots, x_{n})\) （它对 \(\mathrm{M}/(x_{1} \mathrm{~M} + \cdots + x_{i} \mathrm{~M})\) 是正则的）相伴的正合列。记

\[
\theta_ {(x _ {1}, \dots , x _ {i})} \in \operatorname{Ext} _ {\mathbf {A}} ^ {i} (\mathbf {M} / (x _ {1} \mathbf {M} + \dots + x _ {i} \mathbf {M}), \mathbf {M})
\]

\[
\theta_ {(x _ {i + 1}, \dots , x _ {n})} \in \operatorname{Ext} _ {\mathbf {A}} ^ {n - i} (\mathbf {M} / (\boldsymbol {x}) \mathbf {M}, \mathbf {M} / (x _ {1} \mathbf {M} + \dots + x _ {i} \mathbf {M}))
\]

为与 (20) 和 (21) 相伴的扩张类，则根据 X，第 118 页，命题 3，我们有

\[
\theta_ {(x _ {1}, \dots , x _ {n})} = \theta_ {(x _ {1}, \dots , x _ {i})} \circ \theta_ {(x _ {i + 1}, \dots , x _ {n})}.\tag{22}
\]

此外，根据命题 5（X，第 157 页），科斯居尔复形 K.(x, M) 在次数 \(\neq n\) 时零调，由此得到一个正合列

\[
0 \longrightarrow \mathrm{M} \xrightarrow {\partial^ {0}} \mathbf {K} ^ {1} (\boldsymbol {x}, \mathrm{M}) \xrightarrow {\partial^ {1}} \mathbf {K} ^ {2} (\boldsymbol {x}, \mathrm{M}) \xrightarrow {\partial^ {2}} \dots \longrightarrow \mathbf {K} ^ {n} (\boldsymbol {x}, \mathrm{M}) \xrightarrow {q} \mathrm{M} / (\boldsymbol {x}) \mathrm{M} \longrightarrow 0\tag{23}
\]

其中我们已把 \(\mathbf{K}^{0}(\boldsymbol{x},\mathbf{M})\) 与 \(\mathbf{M}\) 等同，且其中 \(q\) 把每个元素 \(\boldsymbol{m}\in\mathbf{K}^{n}(\boldsymbol{x},\mathbf{M})\) 映到 \(\mathbf{M}/(\boldsymbol{x})\mathbf{M}\) 中 \(\boldsymbol{m}(1,2,\ldots,n)\in\mathbf{M}\) 所在的类。

命题 8. --- 假设序列 x 对 M 是正则的。Ext \(_{A}^{n}\) (M/(x) M, M) 中与正合列 (23) 相伴的元素是 (-1) \(^{n(n+1)/2}\) . θ \(_{x}\) 。

对 \(i = 0,1,\dots,n\) ，让我们定义一个 A-线性映射

\[
a ^ {i}: \mathbf {K} ^ {i} (\boldsymbol {x}, \mathrm{M}) \rightarrow \mathrm{M} _ {i} = \mathrm{M} / (x _ {1} \mathrm{M} + \dots + x _ {i - 1} \mathrm{M})
\]

如下：如果 \(\boldsymbol{m} \in \mathbf{K}^i(\boldsymbol{x}, \mathbf{M})\) ，则 \(a^i(\boldsymbol{m})\) 是 \(\mathbf{M}_i\) 中元素 \(\boldsymbol{m}(1, 2, \ldots, i)\) （属于 M）的类。显然， \(a^0\) 是 M 的恒等映射，并且 \(p \circ a^n = q\) 。此外， \(a^{i+1} \circ \partial^i(\boldsymbol{m})\) 是下述元素在 \(\mathbf{M}_{i+1}\) 中的像

\[
\sum_ {k = 1} ^ {i + 1} (- 1) ^ {k + 1} x _ {k} \boldsymbol {m} (1, 2, \dots , k - 1, k + 1, \dots , i + 1).
\]

由于 \(x_{k}\) 零化 \(\mathbf{M}_{i+1}\) 对 \(k=1,\ldots,i,a^{i+1}\circ\partial^{i}(\boldsymbol{m})\) 也是下述元素的像

\[
(- 1) ^ {i} x _ {i + 1} \boldsymbol {m} (1, 2, \dots , i),
\]

因此

\[
a ^ {i + 1} \circ \partial^ {i} = (- 1) ^ {i} \overline {{{x}}} _ {i + 1} \circ a ^ {i}.
\]

根据 X，第 120 页，推论 1 和 2，与 (23) 相伴的 Ext \(_{A}^{n}\) (M/(x) M, M) 的元素等于 \(\prod_{i=1}^{n} (-1)^{i} \cdot \theta_{x}\) ，由此得该断言。

推论。 — 进一步假设诸模 \(\mathbf{M}/(x_{1}\mathbf{M}+\cdots+x_{i-1}\mathbf{M})\) 对 (x)-进拓扑是分离的，并设 \((a_{ij})\in\mathbf{GL}_{n}(\mathbf{A})\) 。令

\[
y _ {i} = \sum_ {j} a _ {i j} x _ {j} \text { and } \mathbf {y} = (y _ {1}, \dots , y _ {n}).
\]

则序列 y 对 M 是正则的，并且有 \(\theta_{y} = \det (a_{ij})^{-1}\theta_{x}\) 。

事实上，序列 y 对 M 是正则的，由命题 6（X，第 158 页）和定理 1（X，第 160 页）可得；最后一个断言由命题 8 以及 X 第 119 页的命题 4 得出。

命题 9. --- 假设序列 x 对 M 是正则的。若 N 是满足 (x) N = 0 的 A-模，则我们有 Ext \(_{A}^{i}\) (N, M) = 0 对 i \textless{} n，并且映射 α → θ \(_{x}\) ∘ α 从 Hom \(_{A}\) (N, M/(x) M) 到 Ext \(_{A}^{n}\) (N, M) （它也是与 (19) 相伴的迭代连接同态，参见 X，第 127 页，推论 3）是双射的。

仍需证明同态 \(\psi^{i}:\alpha\mapsto\theta_{x}\circ\alpha\) 从 \(\operatorname{Ext}_{\mathbf{A}}^{i-n}(\mathbf{N},\mathbf{M}/(\mathbf{x})\mathbf{M})\) 到 \(\operatorname{Ext}_{\mathbf{A}}^{i}(\mathbf{N},\mathbf{M})\) 对 \(i\leqslant n\) 是双射的。让我们对 \(n\) 作归纳推理，该断言对 \(n=0\) 是平凡的。令 \(\mathbf{M}_{1}=\mathbf{M}/x_{1}\mathbf{M}\) , \(\mathbf{x}^{\prime}=(x_{2},...,x_{n})\) ，使得 \(\mathbf{x}^{\prime}\) 对 \(\mathbf{M}_{1}\) 是正则的。由归纳假设，同态

\[
\overline {{{{\psi}}}} ^ {i - 1}: \alpha \mapsto \theta_ {\boldsymbol {x} ^ {\prime}} \circ \alpha
\]

从 \(\operatorname{Ext}_{\mathbf{A}}^{i-n}(\mathbf{N}, \mathbf{M}/(\mathbf{x})\mathbf{M})\) 到 \(\operatorname{Ext}_{\mathbf{A}}^{i-1}(\mathbf{N}, \mathbf{M}_1)\) 对 \(i < n\) 是双射的。此外，考虑正合列

\[
0 \longrightarrow \mathrm{M} \xrightarrow {(x _ {1}) _ {\mathrm{M}}} \mathrm{M} \longrightarrow \mathrm{M} _ {1} \longrightarrow 0;
\]

与之相伴的连接同态 \(\operatorname{Ext}_{\mathbf{A}}^{i}(\mathbf{N}, \mathbf{M}_{1}) \to \operatorname{Ext}_{\mathbf{A}}^{i+1}(\mathbf{N}, \mathbf{M})\) 为 \(\varphi^{i}: \beta \mapsto \theta_{x_{1}} \circ \beta\) （X，第 125 页，命题 5）；

由于我们有 \(\operatorname{Ext}^i(1_{\mathbf{N}}, (x_1)_{\mathbf{M}}) = \operatorname{Ext}^i((x_1)_{\mathbf{N}}, 1_{\mathbf{M}}) = 0\) ，我们推出正合列

\[
0 \longrightarrow \operatorname{Ext} _ {\mathrm{A}} ^ {i} (\mathrm{N}, \mathrm{M}) \longrightarrow \operatorname{Ext} _ {\mathrm{A}} ^ {i} (\mathrm{N}, \mathrm{M} _ {1}) \xrightarrow {\varphi^ {i}} \operatorname{Ext} _ {\mathrm{A}} ^ {i + 1} (\mathrm{N}, \mathrm{M}) \longrightarrow 0.
\]

由于 \(\operatorname{Ext}_{\mathbf{A}}^{i}(\mathbf{N}, \mathbf{M}_{1}) = 0\) 对 \(i < n - 1\) ，我们推出 \(\operatorname{Ext}_{\mathbf{A}}^{i + 1}(\mathbf{N}, \mathbf{M}) = 0\) 对 \(i < n - 1\) ，即 \(i + 1 < n\) ；由此可知 \(\varphi^{i}\) 对 \(i < n\) 是双射的。由于 \(\psi^{i}(\alpha) = \theta_{x} \circ \alpha = \theta_{x_{1}} \circ \theta_{x'} \circ \alpha = \varphi^{i - 1} \circ \overline{\psi}^{i}(\alpha)\) 对 \(\alpha \in \operatorname{Ext}_{\mathbf{A}}^{i - n}(\mathbf{N}, \mathbf{M}/(\mathbf{x})\mathbf{M})\) ， \(\psi^{i}\) 对 \(i \leqslant n\) 确实是双射的。

\backmatter
